\documentclass[14pt,usepdftitle=false,aspectratio=169]{beamer}
\usepackage{preambule}
\setbeamercolor{structure}{fg=black}
\usepackage[nopar]{bigoperators}\usepackage{polynomes,al}\newcommand{\corresp}[8]{\begin{array}[t]{@{}r@{}c@{}l@{}}#1&{}\longleftrightarrow{}&#2\\#4&{}\overset{\scriptscriptstyle#3}{\longmapsto}{}&#5\\#7&{}\overset{\scriptscriptstyle#6}{\longmapsfrom}{}&#8\end{array}}\DeclareMathOperator{\olddom}{dom}\newcommand{\dom}[1]{\olddom\l#1\r}\newcommand{\dashto}{\mathrel{-\mkern-6mu{\to}\mkern-22.5mu{\text{\color{white}\boldmath${\cdot}{\cdot}{\cdot}$}}\mkern1mu\phantom{.}}}
\hypersetup{pdftitle={Introduction à la géométrie algébrique -- Géométrie projective et birationnelle}}
\title{Introduction à la géométrie algébrique\\\emph{Géométrie projective et birationnelle}}
\author{}
\date{}
\begin{document}
\begin{frame}
    \titlepage
\end{frame}
\slideq{Application rationnelle $f\!:\!V\to\bbP^n_\bbk$}{1/12}
\slider{Application partielle $f=\l f_0,\cdots,f_n\r$ avec $f_i\in\fr kV$\linebreak$f$ est régulière en $P$ si chaque $f_i$ le sont et il existe $i$ tel que $f_i\l P\r\neq0$}{1/12}
\slideq{Produit de variétés}{2/12}
\slider{Si $V$ et $W$ sont deux variétés (affines ou projectives) alors $V\times W$ est une variété pour la topologie de Zariski sur le produit ($\bbA_\bbk^m\times\bbA_\bbk^n$ ou $\bbP_\bbk^m\times\bbP_\bbk^n\subset\bbP_\bbk^{\l m+1\r\l n+1\r-1}$ via $\l\lc\underline X\rc,\lc\underline Y\rc\r\mapsto\l X_0Y_0,\cdots,X_0Y_n,\cdots,X_mY_0,X_mY_n\r$)}{2/12}
\slideq{CNS pour que $f\!:\!V\dashto W$ soit birationnelle}{3/12}
\slider{$f$ est dominante ($f\l\dom f\r$ est dense dans $W$) et $f^*\!:\!\fr kW\to\fr kV$ est un isomorphisme\linebreak Il existe $V_0\subset V$ et $W_0\subset W$ deux ouverts non vides tels que $f_{|V_0}^{|W_0}$ est un isomorphisme}{3/12}
\slideq{Recollement affine d'une variété projective}{4/12}
\slider{$V\mapsto V\cap\bbA^n_{\l0\r}=V_{\l0\r}$ et $\overline W\mapsfrom W$ induisent $\set{V\subset\bbP^n_\bbk,V\nsubseteq\l X_0=0\r,\text{irréductible algébrique}}$\linebreak${}\overset{\sim}{\longleftrightarrow}\set{V\subset\bbA^n_{\l0\r}\text{ algébrique irréductible}}$\linebreak$I^a=\set{f\l1{,}X_1{,}{\cdot}{\cdot}{\cdot}{,}X_n\r,f\in I^h}$\linebreak$I^h=\set{X_0^df\l\oldfrac{X_1}{X_0}{,}{\cdot}{\cdot}{\cdot}{,}\oldfrac{X_n}{X_0}\r,f\in I^a}$\linebreak$\fr kV\cong\fr k{V_{\l0\r}}$ et $\dom{f_{V_{\l0\r}}}=\dom f\cap\bbA^n_{\l0\r}$}{4/12}
\slideq{$f\!:\!V\dashto W$ est une application birationnelle}{5/12}
\slider{Il existe $g\!:\!W\dashto V$ telle que $f\circ g=\id_W$ et $g\circ f=\id_V$}{5/12}
\slideq{Morphisme $f\!:\!U\to V$ pour $U\subset V$ un ouvert}{6/12}
\slider{Restriction de $f\!:\!V\dashto W$ où $U\subset\dom f$}{6/12}
\slideq{Nullstellensatz projectif}{7/12}
\slider{Si $\bbk$ est algébriquement clos alors\linebreak$V\l J\r=\varnothing\Leftrightarrow\sqrt J\supseteq\l X_0,\cdots,X_n\r$\linebreak$V\l J\r=\varnothing\Rightarrow I\l V\l J\r\r=\sqrt J$}{7/12}
\slideq{Correspondance $V$--$I$ homogène}{8/12}
\slider{$\corresp{\set{\substack{\textstyle J\leqslant\pol kV\\\textstyle\text{homogènes}}}}{\set{X\subset\bbP^n_\bbk}}V{J}{\set{\substack{\textstyle P\in\bbP^n_\bbk,\\\textstyle\forall f\in J\\\textstyle\text{homogène},\\\textstyle f\l P\r=0}}}I{\set{\substack{\textstyle f\in\pol k{X_0,\cdots,X_n}\\\textstyle\text{homogène},\\\textstyle\forall P\in X,\\\textstyle f\l P\r=0}}}{X}$}{8/12}
\slideq{Bijections induites par $V$ et $I$}{9/12}
\slider{$\set{\l X_0{,}{\cdot}{\cdot}{\cdot}{,}X_n\r\subsetneq J\leqslant\pol k{X_0{,}{\cdot}{\cdot}{\cdot}{,}X_n}\text{ radicaux}}$\linebreak${}\longleftrightarrow\set{X\subset\bbP^n_\bbk\text{ algébriques}}$\linebreak et\linebreak$\set{\l X_0{,}{\cdot}{\cdot}{\cdot}{,}X_n\r\subsetneq J\leqslant\pol k{X_0{,}{\cdot}{\cdot}{\cdot}{,}X_n}\text{ premiers}}$\linebreak${}\longleftrightarrow\set{X\subset\bbP^n_\bbk\text{ algébriques irréductibles}}$}{9/12}
\slideq{Variété rationnelle}{10/12}
\slider{Variété $V$ telle que $\fr kV$ est une extension pûrement transcendante de $\bbk$\linebreak Variété $V$ telle qu'il existe $V_0\subset V$ un ouvert non vide isomorphe à un ouvert non vide $U_0\subset\bbA_\bbk^n$}{10/12}
\slideq{$I$ est un idéal homogène}{11/12}
\slider{Si $f=\sum{i=0}d{f_i}\in I$ avec $f_i$ homogène de degré $i$ alors $f_i\in I$ pour tout $I$}{11/12}
\slideq{Fonction rationnelle sur $V$}{12/12}
\slider{$f\!:\!V\dashto\bbk$ donnée par $f=\frac gh$ avec $g$ et $h$ deux polynômes homogènes de même degré}{12/12}
\end{document}